\documentclass[a4paper]{article}

\usepackage[margin=1in]{geometry}
\usepackage{amsmath,amsfonts,mathtools,amsthm,amssymb}
\usepackage[l2tabu,orthodox]{nag}
\usepackage{microtype} % fixes spacing or whatever
\usepackage{enumitem}
\usepackage{tikz-cd}
\usepackage{xcolor}
\usepackage{physics}
\usepackage{mathrsfs}
\usepackage{cancel}
% \usepackage{libertine,libertinust1math}
\usepackage{eucal}
\usepackage[toc,page]{appendix}
\usepackage{pgfplots}
\pgfplotsset{compat=1.18}

\newtheorem{proposition}{Proposition}
\newtheorem{theorem}{Theorem}
\newtheorem{lemma}{Lemma}
\newtheorem{corollary}{Corollary}

\theoremstyle{definition}
\newtheorem{remark}{Remark}
\newtheorem{definition}{Definition}
\newtheorem{example}{Example}

% \usepackage[sc,osf]{mathpazo}   % With old-style figures and real smallcaps.
% \linespread{1.025}              % Palatino leads a little more leading
% % Euler for math and numbers
% \usepackage[euler-digits,small]{eulervm}
% \AtBeginDocument{\renewcommand{\hbar}{\hslash}}

\usepackage{etoolbox}
\usepackage{dynkin-diagrams}
\def\row#1/#2!{\dynkin{#1}{#2}\\}
\newcommand{\tble}[1]{
   \renewcommand*\do[1]{\row##1!}
   \[
      \begin{array}{ll}\docsvlist{#1}\end{array}
   \]
}

\allowdisplaybreaks

\renewcommand{\epsilon}{\varepsilon}
% \renewcommand{\phi}{\varphi}

\DeclareMathAlphabet{\mathpzc}{OT1}{pzc}{m}{it}

\newcommand{\goth}[1]{\mathfrak{#1}}
\newcommand{\red}[1]{#1_{\text{red}}}
\newcommand{\ad}[1]{#1_{\text{Ad}}}
\newcommand{\reg}[1]{#1_{\text{reg}}}
\newcommand{\sh}[1]{#1^{\text{sh}}}
\newcommand{\cpdv}[2]{\cfrac{\partial #1}{\partial #2}}

\newcommand{\fA}{\mathcal{A}}
\newcommand{\fB}{\mathcal{B}}
\newcommand{\fC}{\mathcal{C}}
\newcommand{\fD}{\mathcal{D}}
\newcommand{\fE}{\mathcal{E}}
\newcommand{\fF}{\mathcal{F}}
\newcommand{\fG}{\mathcal{G}}
\newcommand{\fH}{\mathcal{H}}
\newcommand{\fI}{\mathcal{I}}
\newcommand{\fJ}{\mathcal{J}}
\newcommand{\fK}{\mathcal{K}}
\newcommand{\fL}{\mathcal{L}}
\newcommand{\fM}{\mathcal{M}}
\newcommand{\fN}{\mathcal{N}}
\newcommand{\fO}{\mathcal{O}}
\newcommand{\fP}{\mathcal{P}}
\newcommand{\fQ}{\mathcal{Q}}
\newcommand{\fR}{\mathcal{R}}
\newcommand{\fS}{\mathcal{S}}
\newcommand{\fT}{\mathcal{T}}
\newcommand{\fX}{\mathcal{X}}
\newcommand{\fY}{\mathcal{Y}}
\newcommand{\fZ}{\mathcal{Z}}

\newcommand{\be}{\mathbf{e}}
\newcommand{\bx}{\mathbf{x}}
\newcommand{\bone}{\mathbf{1}}
\newcommand{\cx}{\bar{x}}
\newcommand{\cy}{\bar{y}}
\newcommand{\vx}{\vec{x}}
\newcommand{\vsigma}{\vec{\sigma}}
\newcommand{\tzeta}{\tilde{\zeta}}

\newcommand{\gm}{\goth{m}}
\newcommand{\gp}{\goth{p}}
\newcommand{\gF}{\goth{F}}
\newcommand{\gU}{\goth{U}}
\newcommand{\gV}{\goth{V}}

\newcommand{\A}{\mathbf{A}}
\newcommand{\C}{\mathbf{C}}
\newcommand{\F}{\mathbf{F}}
\newcommand{\G}{\mathbf{G}}
\newcommand{\bH}{\mathbf{H}}
\newcommand{\N}{\mathbf{N}}
\newcommand{\PP}{\mathbf{P}}
\newcommand{\Q}{\mathbf{Q}}
\newcommand{\R}{\mathbf{R}}
\newcommand{\Z}{\mathbf{Z}}

\newcommand{\dlog}{\dd\log}

\newcommand\srestr[2]{{\left.\kern-\nulldelimiterspace #1\vphantom{\small|} \right|_{#2}}}
\newcommand\restr[2]{{\left.\kern-\nulldelimiterspace #1 \vphantom{\big|} \right|_{#2}}}
\newcommand\gsec[2]{\Gamma({#1}, {#2})}
\newcommand\gsecx[1]{\Gamma(X, {#1})}

\DeclareMathOperator{\chr}{char}
\DeclareMathOperator{\rProj}{\mathbf{Proj}}
\DeclareMathOperator{\rSpec}{\mathbf{Spec}}
\DeclareMathOperator{\rHom}{\mathbf{Hom}}
\DeclareMathOperator{\rExt}{\mathbf{Ext}}
\DeclareMathOperator{\id}{id}
\DeclareMathOperator{\Frac}{Frac}
\DeclareMathOperator{\rk}{rank}
\DeclareMathOperator{\deter}{det}
\DeclareMathOperator{\pic}{Pic}
\DeclareMathOperator{\cacl}{CaCl}
\DeclareMathOperator{\trd}{tr.d.}
\DeclareMathOperator{\cl}{Cl}
\DeclareMathOperator{\depth}{depth}
\DeclareMathOperator{\codim}{codim}
\DeclareMathOperator{\Div}{Div}
\DeclareMathOperator{\coker}{coker}
\DeclareMathOperator{\len}{length}
\DeclareMathOperator{\height}{height}
\DeclareMathOperator{\supp}{Supp}
\DeclareMathOperator{\proj}{Proj}
\DeclareMathOperator{\im}{im}
\DeclareMathOperator{\Hom}{Hom}
\DeclareMathOperator{\Der}{Der}
\DeclareMathOperator{\spec}{Spec}
\DeclareMathOperator{\Aut}{Aut}
\DeclareMathOperator{\ch}{char}
\DeclareMathOperator{\tor}{Tor}
\DeclareMathOperator{\Ann}{Ann}
\DeclareMathOperator{\Syl}{Syl}
\DeclareMathOperator{\Sym}{Sym}
\DeclareMathOperator{\GL}{GL}
\DeclareMathOperator{\SL}{SL}
\DeclareMathOperator{\Stab}{Stab}
\DeclareMathOperator{\Perm}{Perm}
\DeclareMathOperator{\Orb}{Orb}
\DeclareMathOperator{\Gal}{Gal}
\DeclareMathOperator{\Supp}{Supp}
\DeclareMathOperator{\Ext}{Ext}
\DeclareMathOperator{\Tor}{Tor}
\DeclareMathOperator{\PGL}{PGL}
\DeclareMathOperator{\hd}{hd}
\DeclareMathOperator{\pd}{pd}
\DeclareMathOperator{\SO}{SO}
\DeclareMathOperator{\SU}{SU}
\DeclareMathOperator{\Sp}{Sp}
\DeclareMathOperator{\Oth}{O}
\DeclareMathOperator{\Uni}{U}
\DeclareMathOperator{\evf}{ev}
\DeclareMathOperator{\cH}{H}
\DeclareMathOperator{\dR}{R}
\DeclareMathOperator{\End}{End}
\DeclareMathOperator{\Hasse}{Hasse}
\DeclareMathOperator{\Num}{Num}

\author{James Lee}

% Solarized
% \pagecolor[RGB]{0,20,26}
% \color[RGB]{191,191,191}

% Sephia
% \pagecolor[RGB]{249,239,220}

% Grey
% \pagecolor[RGB]{200, 200, 200}

% Black
% \pagecolor[RGB]{0,0,0}
% \color[RGB]{255,255,255}


\title{Chapter 4, Section 6}

\begin{document}

\maketitle

\begin{enumerate} [label=\textbf{\arabic*.}, leftmargin=0em]

    \item A rational curve of degree 4 in $\P^3$ is contained in a unique quadric surface $Q$, and $Q$ is necessarily nonsingular.

          \begin{proof}
              Let $X$ be a rational curve of degree 4 in $\P^3$, and let $\fI$ be its ideal sheaf.
              Then we have an exact sequence
              \[ \begin{tikzcd}
                      0 \arrow[r] & \fI \arrow[r] & \fO_{\P^3} \arrow[r] & \fO_X \arrow[r] & 0.
                  \end{tikzcd} \]
              Twisting by $\fO_{\P^3}(2)$ and taking cohomology gives
              \[ \begin{tikzcd}
                      0 \arrow[r] & \cH^0(\P^3, \fI(2)) \arrow[r] & \cH^0(\P^3, \fO_{\P^3}(2)) \arrow[r] & \cH^0(X, \fO_X(2)) \arrow[r] & \cdots.
                  \end{tikzcd} \]
              Now the middle vector space has dimension 10, and $\fO_X(2)$ corresponds to a linear system of degree 8, so it has dimension at most 9 by Riemann-Roch.
              Then, $\cH^0(\P^3, \fI(2))$ has dimension at least 1, so $X$ is contained in a quadric surface $Q$.
              If $X$ is contained in another quadric surface, then it must be a complete intersection by reason of its degree, which implies $\omega_X = \fO_X$ (II, Ex. 8.4), a contradiction.
              If $Q$ is singular, then $X$ must have genus 1 (6.4.1), which is a contradiction.
              Hence, $Q$ must be a nonsingular quadric surface.
          \end{proof}

    \item A rational curve of degree 5 in $\P^3$ is always contained in a cubic surface, but there are such curves which are not contained in any quadric surface.

          \begin{proof}
              Let $X$ be a rational curve of degree 5 in $\P^3$, and let $\fI$ be its ideal sheaf.
              Then we have an exact sequence
              \[ \begin{tikzcd}
                      0 \arrow[r] & \fI \arrow[r] & \fO_{\P^3} \arrow[r] & \fO_X \arrow[r] & 0.
                  \end{tikzcd} \]
              Twisting by $\fO_{\P^3}(3)$ and taking cohomology gives
              \[ \begin{tikzcd}
                      0 \arrow[r] & \cH^0(\P^3, \fI(3)) \arrow[r] & \cH^0(\P^3, \fO_{\P^3}(3)) \arrow[r] & \cH^0(X, \fO_X(3)) \arrow[r] & \cdots.
                  \end{tikzcd} \]
              Now, the middle vector space has dimension 20, and the vector space on the right has dimension 16 by Riemann-Roch on $X$ (note that $\fO_X(3)$ corresponds to a divisor of degree 15).
              So we conclude that $h^0(\P^3, \fI(3)) \geq 4$.

              A degree 5 curve in $\P^3$ contained in a quadric cone must have genus 2 (6.4.1), so if $X$ is contained in a quadric surface, then it must be a nonsingular quadric surface.
              Choosing a degree 5 embedding of $\P^1$ in $\P^3$ is the same as picking a 4-dimensional subspace of $\cH^0(\P^1, \fO_{\P^1}(5))$, which has dimension 6.
              Thus, the family of embeddings $\P^1 \to \P^3$ of degree 5 is parametrized by the Grassmannian $\Gr(4, 6)$, which has dimension $4(6 - 4) = 8$.
              On the other hand, if $X$ is contained in a nonsingular quadric surface $Q \subset \P^3$, then it corresponds to a divisor of type $(1, 4)$ on $Q$.
              So choosing an embedding of $\P^1$ as a curve of type $(1, 4)$ in $Q$ is the same as picking a 2-dimensional subspace of $\cH^0(\P^1, \fO_{\P^1}(1)$ and a 2-dimensional subspace of $\cH^0(\P^1, \fO_{\P^1}(4))$.
              Thus, the family of embeddings $\P^1 \to Q$ is parametrized by $\Gr(2, 2) \times \Gr(2, 5)$, which has dimension $1(2 - 1) + 2(5 - 2) = 7$.
              Since the family of embeddings $Q \to \P^3$ has dimension one, there are plenty of embeddings of $\P^1$ as a degree 5 curve in $\P^3$ which are not contained in any quadric surface.
          \end{proof}

    \item A curve of degree 5 and genus 2 in $\P^3$ is contained in a unique quadric surface $Q$.
          Show that for any abstract curve $X$ of genus 2, there exist embeddings of degree 5 in $\P^3$ for which $Q$ is nonsingular, and there exist other embeddings of degree 5 for which $Q$ is singular.

          \begin{proof}
              Let $X$ be a curve of degree 5 and genus 2 in $\P^3$, and let $\fI$ be its ideal sheaf.
              Then we have an exact sequence
              \[ \begin{tikzcd}
                      0 \arrow[r] & \fI \arrow[r] & \fO_{\P^3} \arrow[r] & \fO_X \arrow[r] & 0.
                  \end{tikzcd} \]
              Twisting by $\fO_{\P^3}(2)$ and taking cohomology gives
              \[ \begin{tikzcd}
                      0 \arrow[r] & \cH^0(\P^3, \fI(2)) \arrow[r] & \cH^0(\P^3, \fO_{\P^3}(2)) \arrow[r] & \cH^0(X, \fO_X(2)) \arrow[r] & \cdots.
                  \end{tikzcd} \]
              The middle vector space has dimension 10, and the vector space on the right has dimension 9 by Riemann-Roch on $X$ (note that $\fO_X(2)$ corresponds to a nonspecial divisor of degree 10).
              So we conclude that $h^0(\P^3, \fI(2)) \geq 1$.
              By reason of degree, $X$ cannot be contained in two distinct quadric surfaces, so there is a unique irreducible quadric surface $Q$ containing $X$.

              Similar to (Ex. 6.2), if $X$ is a genus 2 curve of degree 5 in a nonsingular quadric surface $Q \subset \P^3$, then it corresponds to a divisor of type $(2, 3)$ on $Q$.
              Thus, choosing an embedding of $X$ as a curve of type $(2, 3)$ in $Q \subset \P^3$ is the same as picking a $g_2^1$ and $g_3^1$ on $X$.
              Since $X$ is hyperelliptic, it always has a $g_2^1$ and $g_3^1$, so there are plenty of embeddings of $X$ into a nonsingular quadric surface.
              In fact, any $g_3^1$ on $X$ is obtained by adding a point to the unique $g_2^1$, so the family of $g_3^1$'s on $X$ has dimension one, implying that the family of embeddings of $X$ into a nonsingular quadric has dimension one.
              On the other hand, any effective divisor of degree 5 gives an embedding of $X$ in $\P^3$ (Ex. 3.1), which is the same as picking 5 points in $X$, so the family of embeddings $X \to \P^3$ has dimension at least 5.

          \end{proof}

    \item There is no curve of degree 9 and genus 11 in $\P^3$.

          \begin{proof}
              Suppose there exists a curve $X \subseteq \P^3$ of degree 9 and genus 11, and let $\fI$ be its ideal sheaf.
              We have an exact sequence
              \[ \begin{tikzcd}
                      0 \arrow[r] & \cH^0(\P^3, \fI(2)) \arrow[r] & \cH^0(\P^3, \fO_{\P^3}(2)) \arrow[r] & \cH^0(X, \fO_X(2)) \arrow[r] & \cdots.
                  \end{tikzcd} \]
              Then $\fO_X(2)$ is a very ample line bundle corresponding to a linear system of degree 18, say $|\fO_X(2)| = |D|$ for some very ample divisor $D$ of degree 18.
              Riemann-Roch gives
              \begin{equation*}
                  \dim |D| - \dim |K - D| = 8.
              \end{equation*}
              If $D$ is nonspecial, then $\dim|D| = 8$.
              If $D$ is special, Clifford's theorem (5.4) says
              \begin{equation*}
                  \dim|D| \leq \frac{1}{2} \deg{D} = 9,
              \end{equation*}
              and we have an equality if and only if $X$ is hyperelliptic and $D$ is a multiple of the unique $g_2^1$ on $X$.
              But $D$ is very ample, so it cannot be a multiple of the $g_2^1$, and the inequality must be strict.
              We conclude that
              \begin{equation*}
                  h^0(X, \fO_X(2)) = 1 + \dim|D| \leq 9,
              \end{equation*}
              and therefore, $\cH^0(\P^3, \fI(2)) \neq 0$.
              Hence, $X$ must be contained in a quadric surface.

              Following the discussion in (6.4.1), if $X$ is contained in nonsingular quadric surface, then it corresponds to a divisor of type $(a, b)$ such that
              \begin{equation*}
                  a + b = 9, \quad ab - a - b + 1 = 11,
              \end{equation*}
              which is not possible.
              If $X$ is contained in a singular quadric surface, then $X$ must have genus 12.
              Hence, there is no curve of degree 9 and genus 11 in $\P^3$.
          \end{proof}

    \item If $X$ is a complete intersection of surfaces of degree $a, b$ in $\P^3$, then $X$ does not lie on any surface of degree $< \min(a, b)$.

          \begin{proof}
              If $X$ is a complete intersection of surfaces $H, H'$ of degree $a, b$ respectively, in $\P^3$, then $X$ has degree $ab$.
              Without loss of generality, suppose $a \leq b$.
              If $H''$ is surface of degree $c < a$ containing $X$, then the irreducible components of $H \cap H''$ can have degree at most $ac$ (I, 7.7), which is strictly less than $ab$, so $X$ cannot be an irreducible component of $H \cap H''$.
              Hence, $X$ cannot lie on any surface of degree $< a$.
          \end{proof}

\end{enumerate}

\end{document}
